\subsubsection{Finite projections for the evaluation at the prime assigned to the sequence}
\label{subsec:word-projection}

The measure description turns the required finite projections into support conditions. We compute them in the Hom lattice to retain the distinction between the distinguished and conjugate places. The complex argument includes the factor
\(\delta\) from the uniformization. Write \(e=e^\natural\)
for the normalized vector, \(\pi\bar\pi=p\), and
\(u=\pi/\bar\pi\). The local row conditions are at the
places dividing \(\pi\). For the moment measure, the
split factor chosen by the coefficient embedding is the
one where \(\rho\) has positive valuation, where
\(\rho=\pi\) in the distinguished branch and
\(\rho=\bar\pi\) in the conjugate branch. The word
generator belongs to the common CM subfield for the
selected embeddings, so all derivation positions use this
same choice.

\begin{proposition}[Moment evaluations and the finite vacuum vector]\label{prop:word-reader-vacuum}
On the marked Hom model, the translates \textup{\eqref{eq:tag-m6}} use the
same moment measure at each selected coefficient place, with the point evaluations \textup{\eqref{eq:tag-m7}} for any finite list of jets.
For \(x\) in the prescribed support, the finite vacuum in \textup{\eqref{eq:tag-m8}} is
\[
 \Theta_x=\frac1n W_{-t}\sum_{z\in R}W_z e,\qquad n=[M\cap A:M\cap A\cap B],
\]
where \(R\) represents \(B/(M\cap B)\). The expression is independent of the
choice of the allowed rational lift \(t\).
\end{proposition}

\begin{proof}
\proofstep{Read the connected torsion on the chosen split factor.}

A translate by
\[
                        y=N_0(\bar\rho/\rho)m,\qquad m\in M,
\tag{M6}\label{eq:tag-m6}
\]
qualifies for the preceding primary calculation with
\(n=1\): \(\bar\rho/\rho=\bar\rho^2/p\) and
the minus component of \(\bar\rho^2\) is divisible by
\(p^2\). This also applies to the changed companion
\(M''\) below. The prime-to-\(p\) translate may be
written with \(W\) in place of \(U\) at the cost of a
root of unity. Its exponent is rational because prime-to-\(p\)
multiples lie in the lattice and the Hermitian pairing
there is integral over the CM order. Thus whenever
a list of jets is taken after a family of translates
\eqref{eq:tag-m6}, the moment calculation gives simultaneous
approximations on \(M\) with point evaluations
\[
                      \lambda(\kappa(z,y))P_m(z),
\tag{M7}\label{eq:tag-m7}
\]
where \(P_m(z)\) denotes the monomial of the adjoint
entries with the units in \eqref{eq:tag-m5}. All jets are algebraic
in these frames. The same coefficients approximate
every index in the list, including when the list
contains different \(y\)'s.

\proofstep{Compute the two half lattices and their integral cocycles.}

Set

\[
  S=S_G\cup\{P:P\mid\pi\},
\]
 where \(S_G\) is the set
of away positions and is disjoint from its conjugate.
Write \(H_S\) for the ordinary Hom conditions at \(S\)
with no conditions elsewhere, and let \(c_0,g\) be the
prime-to-\(p\) denominator and level factors. At every
place of \(S\) or its conjugate the period coefficient
ideal is a unit, so \(M\) there is the Hom lattice of
the selected self-dual source and target. Away from these
places the coefficient ideal stays in \(M\).
The half lattice \(B\) is specified place by place by
\[
  B_P=\begin{cases}pg\,M_P&P\in S,\\ (pg)^{-1}M_P&P\in\bar S,\\ M_P&P\notin S\cup\bar S.\end{cases}
\]
The other half lattice \(A\) uses the reciprocal scalings for \(c_0pg\).
Both are integral for the Weyl cocycle.

In fact the
product of the scaling ideals with their conjugates is
the unit ideal, so the Hermitian integrality
\(\beta(M,M)\subset O\) also holds on each of these
lattices. For \(v,w\) in either half lattice, take \(a,b\in\mathbb Z\)
with 
\[
  \beta(v,w)=a+b(1+\sqrt{-d})/2.
\]
 The rational cocycle in
\(W_v W_w\) is then integral by the identity
\[
             \RePart\beta(v,w)+\tfrac12\operatorname{Im}H(v,w)=a+b.
\]

\proofstep{Lift the marked indicator in its actual support.}

For \(x\in c_0^{-1}H_S\) contributing to the reader,
the local marked indicator imposes
\(xD_v\subset D'_v\) and \(\ell'_v x=\ell_v\)
at all \(v\mid\pi\). Choose a rational lift \(t\in A\)
with \(t-x\in pgH_S\). Such a lift exists by
approximation in the cosets at \(S\), with the remaining
integral conditions on \(A\). At \(P\mid\pi\) this lift is
integral because \(x\) is integral and the congruence
has modulus divisible by \(p\). At the opposite primes
it is integral because \(A\) has the positive scaling.
It follows that \(c_0gt\in M\) at every place. From the
definition of \(B\), \(g uM\subset B\), \(g u^{-1}B\subset M\)
and the localization of \(B\) at \(p\) is that of \(uM\).

\proofstep{Evaluate the finite transform and check independence of the lift.}

The moved theta indicator in the reader is
\[
             \Theta_x=\frac1n W_{-t}\sum_{z\in R}W_z e ,
\qquad
        n=[M\cap A:M\cap A\cap B],\quad
        R\text{ represents }B/(M\cap B).
\tag{M8}\label{eq:tag-m8}
\]
In the finite transform of the zero-coset indicator, among
the representatives of \(M\cap A\) modulo \(B\) only the
zero coset contributes. The pairing with \(z\in B\) is
trivial on its representative, by integrality on \(B\).
Thus the sum in the transform is one and the character
average supplies \(1/n\).
The finite transform contributes the square root of the
Hom volume, which cancels the inverse square root in
the theta evaluation.

The Weyl translates
are well defined modulo \(M\cap B\), since these
translations fix \(e\) and the cocycle on \(B\) is
integral. Two lifts have difference in \(B\) by their
congruences and the reciprocal scalings off \(S\).
This difference fixes the sum in \eqref{eq:tag-m8}, and its
composition with the lift contributes no phase by
integrality on \(A\). Thus \eqref{eq:tag-m8} uses the original
indicator for any such lift. In particular
\(n\mid[M:M\cap B]\).
\label{cor:word-projection-1}

\end{proof}

\begin{proposition}[The two projected moment approximations]\label{prop:word-projection-moments}
On the selected pair with the ordinary measure of
Proposition~\ref{prop:theta-measure-word}, the two finite
projections and the away reader give precisely the moment
condition \eqref{eq:tag-w6} used in Proposition~\ref{prop:word-moment-criterion}.
For each local component \(v\) at \(p=T\), the two supports are
\[
 \begin{aligned}
 x^*(\ker\ell'_v)&\subset D_v
   &&\text{at the distinguished place, with factor }p^{-F},\\
 xD_v&\subset D'_v,\quad \ell'_v x|_{D_v}=\ell_v|_{D_v}
   &&\text{at the conjugate place, without that factor}.
 \end{aligned}
\]
The two
denominator powers of \(p\) are uniform in the jet list
and the precision, also after the away reader.
\end{proposition}

\begin{proof}
Use Proposition~\ref{prop:word-reader-vacuum} on these lattices.
Since \(B\) agrees with \(uM\) at \(p\), the
\(p\)-primary part of \(B/(M\cap B)\) has representatives
in \(N_1uM\) for every sufficiently divisible
prime-to-\(p\) integer \(N_1\) with \(g\mid N_1\).
The other part has representatives in \(pB\):
multiplication by \(p\) kills the first part and is
invertible on the second. These statements follow
also by localizing the quotient at \(p\) and away from
\(p\), using \(p u M\subset M\),
\(pM\subset uM\) and \(gpB\subset M\).

Take sets of representatives \(R_1,R_2\) in those
two lattices. Then \(R=R_1+R_2\) in \eqref{eq:tag-m8}, and
\[
        |R_1|=[uM:M\cap uM]=[M:M\cap uM].
\tag{M9}\label{eq:tag-m9}
\]
Both sides of the second equality are the orders of
dual quotients of the self-dual lattices \(M,uM\)
under the alternating pairing. Their common order
has \(p\)-valuation at least \(v_p(n)\), by the
divisibility after \eqref{eq:tag-m8} and \(B_p=(uM)_p\).

\proofstep{Project at the conjugate coefficient place using prime-to-\(p\) translates.}
For each \(z_2\in R_2\), \(z_2-t\) is a prime-to-\(p\)
translate, since 
\[
  gpB\subset M
\]
 and \(c_0gt\in M\).
Take \(N_1\) divisible also by the finitely many
integers \(N_0\) for these translates and by the
denominator of \(t\). All are prime to \(p\).
Every \(z_1\in R_1\) then has the form \eqref{eq:tag-m6}
for \(\rho=\bar\pi\), with \(m\in M\) for each
of these measures. Commuting \(W_{-t}\) past
\(W_{z_1}\) in \eqref{eq:tag-m8} gives the factor
\(\kappa(-t,z_1)\). Combining it with \eqref{eq:tag-m7}
gives \(\kappa(z-t,z_1)\) on the point sample \(z\).
This combination uses multiplicativity of the character in its first variable,
\[
  \kappa(z,z_1)\kappa(-t,z_1)=\kappa(z-t,z_1).
\]
The additional composition of \(W_{-t}\) with
\(W_{z_2}\) gives a root of unity depending on
\(z_2,t\). The cocycle on \(B\) is integral.

For a sample \(z\in M\) the characters of \(z_1\)
descend to the \(p\)-quotient in \eqref{eq:tag-m9}. At \(p\),
\(z-t\) is integral in \(M\), and away from \(p\)
the multiplier \(N_1\) clears the denominator
of \(t\). Thus the pairing with the denominator
lattice of that quotient is integral. Let \((uM)_p^\#\) denote the alternating dual of the local
period lattice. Orthogonality on the finite quotient in \eqref{eq:tag-m9}
gives
\[
 \sum_{z_1\in R_1}\kappa(z-t,z_1)
   = |R_1|\mathbf 1_{(uM)_p^\#}(z-t)
   = |R_1|\mathbf 1_{(uM)_p}(z-t).
\]
The prime-to-\(p\) conditions hold by the choice of \(N_1\).
At \(p\), \(N_1\) is a unit and \(u\bar u=1\) makes \(uM\)
self-dual, giving the last equality.

At each
\(P\mid\pi\) this gives
\(z-t\in pM_P\), and \(t-x\in pgH_S\) gives
\(z-x\in pM_P\). The point sample therefore obeys
the second line of \eqref{eq:tag-w6}. For each surviving sample, write \(c\) for its original coefficient and \(\zeta\) for the accompanying root of unity. Its new coefficient is
\[
  c'=\zeta\,\frac{|R_1|}{n}\,c.
\]
 Equation \eqref{eq:tag-m9} proves that this
coefficient is integral. The simultaneous approximations may be taken as
precise as needed for the finitely many scalars before
summation.
\label{cor:word-projection-2}

\proofstep{Keep the row-kernel support when projecting at the distinguished coefficient place.}

Replace the target with the self-dual companion
\(N''\) whose local lattices at \(P\mid\pi\) are the
target row kernels and whose opposite lattices are
their duals. At every other place leave the target
lattice fixed. These local conditions define a
global rational lattice by intersection. In the
changed Hom period lattice \(M''\), put
\(I=M\cap M''\) and 
\[
  k_0=[M:I]=[M'':I].
\]
The local Hom descriptions give
\[
     \pi M\subset M'',\qquad \bar\pi M''\subset M,\qquad
                         k_0\mid p^{fN_i}.
\tag{M10}\label{eq:tag-m10}
\]
At a first place the changed target imposes one
residue row equation per source coordinate. There
are \(f\) such places, each of residue cardinality
\(p\). At an opposite place \(M''\) enlarges \(M\).
At other primes the lattices agree. The half different
is a unit at these unramified split factors, so the
index is computed on these Hom lattices directly.

Let 
\[
  C=c_0g\bar\pi I
\]
 and let \(J\) represent
\(C/(C\cap pgH_S)\). At the row places \(C\) sends the
source into the row kernel. At the away places
\(C\subset pgH_{S_G}\), because the extra \(p\) is
a unit there and the factor \(g\) clears the level.
Thus both the row indicator and the away reader
are constant on translates by \(C\). Moreover
\(C\subset M\cap A\).

Localizing the condition
\(c_0g\bar\pi m\in pgH_S\) for \(m\in I\) gives
\(m\in\pi M\): at \(P\mid\pi\) it is precisely
the factor of \(p\) in the congruence, and at
\(S_G\) the condition has already been supplied
by \(g\). The other localizations impose no reader
condition and \(\pi\) is a unit there. Therefore
\[
                  |J|=[I:\pi M],\qquad
       v_p(n)\le v_p([M:\pi M])
                  =v_p(k_0)+v_p(|J|).
\tag{M11}\label{eq:tag-m11}
\]
For the inequality use \(g\pi M\subset B\), which
follows from \(guM\subset B\) by multiplying by
\(\bar\pi\). The factor \(g\) has zero \(p\)-valuation.

We may group the total jet sum using \(J\) before
applying a bound. In \(\sum_{j\in J}\Theta_{x+j}\)
the lift in \eqref{eq:tag-m8} is \(t+j\), since it is in \(A\)
and gives the required congruence. Write

\[
  W_{-(t+j)}=W_{-t}W_{-j}
\]
 by integrality on \(A\).
As \(W_{-j}e=e\) for \(j\in C\), commuting it across
\(W_z\) multiplies each term by the pairing character.
It descends to the quotient defining \(J\), since
\(C\cap pgH_S\subset B\) and the pairing is integral
on \(B\). Precisely the \(z\in B\) that pair
integrally with \(C\) survive, with multiplier \(|J|\).

The surviving quotient of \(B\) modulo \(M\) is
the direct sum of the images of \(B\cap M''\) and
\(pB\). To verify it at \(p\), write the paired
local lattices over \(\pi,\bar\pi\) with subscripts
\(+,-\) for this calculation. Then
\[
 \begin{array}{c|cc}
      &+&-\\ \hline
 B   &pM_+&p^{-1}M_-\\
 I   &M''_+&M_-\\
 C   &M''_+&pM_-
 \end{array}
\]
where multiplication by a unit has no effect. The
annihilator of \(C\) in \(B\) is
\(pM_+\oplus M''_-\), by duality between the
two factors. This is \(B\cap M''\) at \(p\).
Away from \(p\) the pairing with \(C\) is
automatically integral, since \(gB\) is in \(M\)
there and \(C\) has the factor \(c_0g\).
The whole quotient away from \(p\) is supplied
by \(pB\). The first summand has image of
order \(k_0\) and represents \(M''/I\).
\label{cor:word-projection-3}

The normalized section for \(M''\) is the image
of the normalized section for \(M\) under the
same theta-line transfer with the determinant
character. Thus, for \(R_1\) representing the
image of \(B\cap M''\),
\[
                       \sum_{z\in R_1} W_z e
                         =\sqrt{k_0}\,\kappa_0 e'',
\tag{M12}\label{eq:tag-m12}
\]
where \(\kappa_0\) is the relative determinant
character to the source half-rank. At each of
the \(f\) places the determinant changes by a
single row index and the opposite determinant
by its inverse. The relative determinant ideal
is generated by \(\pi/\bar\pi\) with the inverse
orientation in the character. Applying the
reflex character gives the product of
\(\tau(\bar\pi/\pi)\) over the \(f\) positions.

Since \(\pi\bar\pi=p\) and each factor before
inversion is integral at \(p\), the norm of a
ratio is at most \(p\). Thus
\(\abs{\lambda(\kappa_0)}_p^2\le p^{fN_i}\).
In \eqref{eq:tag-m12} the section and determinant character
are those of the same basic theta move as in
Section~\ref{subsec:word-measure}.
\label{cor:word-projection-4}

For each representative \(z_2\in pB\) the
remaining term is the translate \(W_{z_2-t}e''\)
times a root of unity and
\[
                             \frac{|J|\sqrt{k_0}}{n}\kappa_0 .
\tag{M13}\label{eq:tag-m13}
\]
The composition of the surviving two sums
gives \eqref{eq:tag-m12} before applying \(W_{-t}\),
since the cocycle is integral on \(B\).
Write \(v=z_2-t\). Some \(q\) prime to \(p\)
has \(qv\in M\), so \(q\pi v\in M''\).
Choose coefficients over the CM integer ring
with \(a q\bar\pi+b\pi=1\). Then
\(v_0=b\pi v\) is a prime-to-\(p\) translate
in \(M''\) and the remaining displacement is
\[
 v-v_0=(1-b\pi)v=aq\bar\pi v=(\bar\pi/\pi)m,
 \qquad m=aq\pi v\in M'' .
\]

For the measure of
\(v_0\), split \(m=N_0 m_1+p m_2\) by
B\'ezout since \(N_0\) is prime to \(p\).
The term \(\bar\pi^2m_2\) belongs to \(M''\).
Therefore \(W_v e''\) is a root of unity times
the primary translate \eqref{eq:tag-m6} of \(W_{v_0}e''\),
by the Weyl law and the fixed lattice of \(e''\).
This allows the moment approximation on \(M''\)
with integral measure coefficients. Each such point
sample has its adjoint carrying the target row
kernel into the source at every \(P\mid\pi\),
by duality with the opposite target lattice.
This is the first line of \eqref{eq:tag-w6}.
By \eqref{eq:tag-m10} and \eqref{eq:tag-m11}, the valuation of \eqref{eq:tag-m13} is at least
\[
                   -\frac{v_p(k_0)}2-\frac{fN_i}{2}
                   \ \ge -fN_i .
\]
The estimate uses the count of the refined reader
cosets before bounding the coefficient. This
recovers the stated power \(p^{-fN_i}\).

\proofstep{Assemble the same finite list at the required precision after applying the away reader.}

The finite jet sum is over
\(c_0^{-1}H_S/(pgH_S)\). Writing \(E_1,E_2\) for the two prescribed row conditions and \(a_{\rm away}(x)\) for the away value, the coefficient at the coset of \(x\) is
\[
  \mathbf1_{E_1}(x)\,\mathbf1_{E_2}(x)\,a_{\rm away}(x).
\]
 In the distinguished
case group the sum by translates of \(C\) before
using \eqref{eq:tag-m13}. In the conjugate case take the
individual cosets. A nonzero coefficient
forces the two row conditions on the original
coset representative \(x\). The away value is
integral at the chosen place and is constant
on the cosets just grouped. The product with
the inverse of the determinant scalar of the
unmodified source and target still has norm at
most one.

For rational jet indices this product
of the scalar with each analytic jet is the
specified coefficient-field jet, under the
chosen coefficient embedding. Since the list
of indices and the number of cosets are finite,
we can choose the approximations before the
finite assembly with precision high enough
for all its scalars.
For instance, given required precisions \(n_i\) at finitely many indices in the same local norm, one may take
\[
  n=\max_i n_i,\qquad p^{-n}\le p^{-n_i}\quad\text{for every }i .
\]
 The ultrametric inequality
preserves integrality at repeated samples.
Consequently the asserted simultaneous
approximations in \eqref{eq:tag-w6} hold for the actual
away reader and coefficient embeddings. The
support and both powers of \(p\) are uniform
in the derivative degrees, precision, reader
and auxiliary coset levels.
\end{proof}

